% Einheit 2 von 3 – Schnittpunkt zweier Geraden (bank/schnittmengen/e2.jsonl, zusammenbau v0.1)
%% TODO Zeitmarke Einheit 2: Marken-Zeile nicht lesbar
%% TODO Zweigzeile Teil 1: was hier gelernt wird (Satz in Schülersprache aus dem Einheitstitel) – Einheit 2
\einheitenkopf[e2]{Einheit 2 von 3 · Schnittpunkt zweier Geraden}
\zweigzeile{Abitur GK}
\setcounter{aufgabe}{12}

%% TODO Titel: Ich-kann-Satz fehlt in der Bank (Platzhalter: Kettenname) – Nr. 13
%% TODO Anweisung: ein Satz über der ersten Teilaufgabe fehlt in der Bank – Nr. 13
\begin{aufgabe}{Schnittpunkt zweier Geraden}
%% TODO \rechenplatz{n}: Schrittzahl des Grundfalls fehlt in der Bank (nur bei fester Schreibform, 2.3 b)
\begin{teile}
\teil Zwei Geraden im Raum werden gleichgesetzt: $(1|0|2) + r \cdot (1|2|0) = (3|4|1) + s \cdot (0|1|1)$. Was entsteht? \\ \kreuz{drei Gleichungen mit zwei Unbekannten} \\ \kreuz{zwei Gleichungen mit zwei Unbekannten} \\ \kreuz{drei Gleichungen mit drei Unbekannten}
\teil Bestimme den Schnittpunkt der Geraden $g: \vec{x} = (2|2|3) + r \cdot (1|0|2)$ und $h: \vec{x} = (3|0|3) + s \cdot (0|1|1)$. S( \leerfeld | \leerfeld | \leerfeld )
\teil Gegeben sind $g: \vec{x} = (1|4|-2) + s \cdot (2|0|3)$ und die Gerade $h$ durch $A(3|0|1)$ und $B(4|1|b)$ mit einer reellen Zahl $b$. $g$ und $h$ haben einen gemeinsamen Punkt. Bestimme $b$. \hfill (Abitur 2023 GK) \\ b = \leerfeld
\teil Zwischen zwei senkrechten Pfählen einer Kletteranlage wird ein Netz gespannt; seine untere Kante verläuft von $(0|0|1{,}5)$ am ersten Pfahl bis zum Punkt $U(6|4|z_U)$ am zweiten Pfahl. Sie berührt die Kante $RT$ einer Plattform mit $R(2|3|2)$ und $T(4|1|2)$; eine Längeneinheit entspricht 1 m. Bestimme die Höhe $z_U$. \hfill (Abitur 2018 LK) \\ $z_U$ = \leerfeld
\teil Ein Museum wird durch einen Körper modelliert: der untere Teil $ABCDEF$ ist Teil der Pyramide $DEFS$, das Dreieck $ABC$ liegt in der $xy$-Ebene, das Dreieck $DEF$ parallel dazu; $A(-1|2|0)$, $B(-1|6|0)$, $D(1|0|6)$, $E(1|8|6)$. Abgedruckt ist die Rechnung \rechnung{(1|0|6) + r \cdot (-2|2|-6) &= (1|8|6) + s \cdot (-2|-2|-6) \\ r = s &= 2 \\ (1|0|6) + 2 \cdot (-2|2|-6) &= (-3|4|-6)} Erläutere das dargestellte Vorgehen. \hfill (Abitur 2018 LK)
\teil Gegeben ist der Würfel $ABCDEFGH$ mit $A(0|0|0)$, $G(4|4|4)$ und $H(0|4|4)$. Für jede reelle Zahl $r$ ist die Gerade $g_r: \vec{x} = (1 - r|0|r^2) + u \cdot (2|4|0)$, $u \in \mathbb{R}$, gegeben. Für einen Wert von $r$ schneidet $g_r$ die Kante $GH$. Bestimme das Verhältnis, in dem der Schnittpunkt die Kante $GH$ teilt. \hfill (Abitur 2019 GK) \\ HQ : QG = \leerfeld : \leerfeld
\end{teile}
\end{aufgabe}

%% TODO Titel: Ich-kann-Satz fehlt in der Bank (Platzhalter: Kettenname) – Nr. 14
%% TODO Anweisung: ein Satz über der ersten Teilaufgabe fehlt in der Bank – Nr. 14
\begin{aufgabe}{Schattenpunkt auf einer Kante als Schnittpunkt von Lichtgerade und Kantengerade berechnen}
\begin{teile}
\teil Ein Spielturm hat ein Dach in Form einer Pyramide; an der Spitze ist eine senkrechte Stange mit dem oberen Ende $T(0|0|7)$ befestigt; der Boden ist die $xy$-Ebene, eine Längeneinheit entspricht 1 m. Sonnenlicht fällt in parallelen Geraden mit dem Richtungsvektor $(4|-1|-2)$. Der Schatten des oberen Endes fällt auf die Dachkante $EF$ mit $E(6|-3|4)$ und $F(6|3|4)$. Bestimme diesen Schattenpunkt. \hfill (Abitur 2017 GK) \\ ( \leerfeld | \leerfeld | \leerfeld )
\end{teile}
\end{aufgabe}

%% TODO Titel: Ich-kann-Satz fehlt in der Bank (Platzhalter: Kettenname) – Nr. 15
%% TODO Anweisung: ein Satz über der ersten Teilaufgabe fehlt in der Bank – Nr. 15
\begin{aufgabe}{Schnittpunkt zweier Geraden – Fehler finden, Begründen}
\begin{teile}
\teil Ich finde den Fehler: Kai prüft, ob sich $g: \vec{x} = (1|0|2) + t \cdot (1|1|0)$ und $h: \vec{x} = (2|2|1) + t \cdot (1|0|1)$ schneiden: \rechnung{1 + t &= 2 + t \\ 1 &= 2} und schreibt: „Kein Schnittpunkt.“
\teil Begründe, warum man beim Gleichsetzen zweier Geraden im Raum nur zwei der drei Koordinatengleichungen zum Lösen braucht, die dritte aber trotzdem prüfen muss.
\end{teile}
\end{aufgabe}
